\section*{Part IX: A Founder's Operating Method}

\subsection*{92. The First Exercise: Trace One Real Unit of Value}

Choose a customer whose transaction has been completed or a pilot that can be run soon. Reconstruct the entire service from 1st inquiry to outcome verification. Record work performed by founders, employees, contractors, software, customers, and external partners. Estimate minutes, error rates, cash payments, and required capabilities.

Do not begin with a corporate chart of accounts. Begin with \textbf{the work actually done}. If there are no customers yet, create a provisional process map and mark every estimate as an assumption. The next pilot should measure which estimates are badly wrong.

\subsection*{93. Classify Every Substantial Activity by Its Cost Behavior}

For each activity ask whether it is fixed for this cohort, variable by learner, triggered by a capacity threshold, or potentially reusable. Record whether its labor requires rare judgment or routine execution. Estimate the setup cost of changing the process and the consequences of failure.

The most actionable findings often look unremarkable: a standard application form eliminates 4 repeated calls; an assignment template reduces technical-support incidents; a weekly group debrief resolves a common misconception that previously required 6 one-on-one sessions.

\subsection*{94. Find the Constraint Before Optimizing Anything}

Mark each stage with demand, available capacity, quality standard, average and peak load, and current waiting. What happens if volume increases by 25\% next month? Which stage fails 1st? What is the cost of 1 extra unit of capacity at that stage?

A founder might discover that lesson production is plentiful, but project review cannot scale. They should then stop investing in more videos and investigate better review workflow, task design, trained raters, and bounded cohort size. Optimization away from the bottleneck is often a sophisticated form of procrastination.

\subsection*{95. Diagnose Repeated Work Before Automating It}

Collect 20 or 30 examples of a recurring task if feasible, appropriately handling privacy and client permission. How much of it is mechanical, patterned, ambiguous, or judgment-intensive? Which cases are rare but costly? What quality standard would a tool need to satisfy? What happens when it is wrong?

Before purchasing an AI system, test a simpler intervention: a checklist, rubric, shared example, standard template, better scheduling, or training. Technology is valuable when it solves a measured problem. The cheapest automation may be \textbf{removing an unnecessary step entirely}.

\subsection*{96. Pilot Changes Using Paired Economic and Quality Measures}

For an automation or process improvement, collect: task completion time, total human review time, direct software expense, setup and maintenance hours, errors missed, necessary rework, escalation rate, customer result, and effect on the next activity. Compare similar work where possible and report uncertainty.

Do not call an intervention successful because it reduced keystrokes or response time. If its purpose is faster feedback, demonstrate that useful, accurate feedback actually arrived sooner and that learner performance did not deteriorate. If its purpose is lower cost, show the cost where it appears in the full process.

\subsection*{97. Define Economic and Educational Stop Conditions}

A pilot should have explicit boundaries. Examples: cumulative cash loss cannot exceed a founder-approved affordable-loss limit; the founder cannot sustain more than a defined weekly workload; essential feedback must remain within a promised window; the assessment standard will not be relaxed simply to improve completion; customer data may not be processed without necessary rights and safeguards.

These are \textbf{decision rules}, not aspirational values. A company that sells more seats after its service-level breaches are already visible is not demonstrating ambition. It is converting future trust into current revenue.

\subsection*{98. A Weekly Operating Review Should Produce Decisions}

A small institution need not hold endless management meetings. 1 weekly review can examine: qualified pipeline, capacity used, upcoming submission peaks, overdue work, quality incidents, learner or client progress, founder-hour utilization, collection status, and the single most important process correction.

Every meeting should end with named responsibility and a date for evaluating the change. A dashboard without a feedback loop is an expensive description of the past.

\subsection*{99. The Monthly Economics Review}

Reconcile actual collections with delivery liabilities; compare realized revenue and variable costs; update founder and staff hours; analyze cost by cohort or contract; inspect overdue receivables and refund exposure; identify step-cost thresholds; and ask whether quality improved, degraded, or remains uncertain.

Also track \textbf{where the model was wrong}. If the founder estimated 10 hours of review but spent 25, the right response is to update pricing and service design: not hide the time because the cash account looks favorable.

\subsection*{100. Build a Capacity Investment Memo Before Hiring or Software Development}

For any material new commitment, write: the constraint being addressed, how it is measured, the present and forecast load, candidate solutions, expected impact on outcomes, setup and recurring costs, quality failure risks, required permissions, cash at risk, reversibility, and stopping or revision triggers.

This forces the founder to distinguish \textbf{capacity investments} from vanity expenditures. A glamorous studio, bespoke platform, or high-status hire may eventually be appropriate. It must earn its place by improving a necessary capability or an economically justified customer experience.

\subsection*{101. When to Reduce Price, Add Volume, or Move Upmarket}

Low demand can tempt a founder to discount tuition until the class fills. But price changes also alter customer composition, commitment, service expectations, and unit contribution. The right action depends on the diagnosis.

If the product is useful but acquisition is expensive, improve buyer targeting and evidence before cutting delivery quality. If the price is high relative to the customer's value, redesign the scope or buyer segment. If qualified demand exceeds reviewer capacity, consider a higher fee, smaller groups, more reliable staff, or different service tiers with genuinely different obligations. If a premium customer needs far more individualized work, price that work instead of subsidizing it invisibly.

Pricing is revisited in Part IV. Here the central lesson is that \textbf{unit economics, service scope, and target customer must be redesigned together}.

\subsection*{102. How to Know When Standardization Has Gone Too Far}

Watch for evidence that the system cannot handle valid differences: advanced learners find tasks trivial, novices receive irrelevant responses, unusual but legitimate questions are rejected, feedback loses specificity, or a procedural rule takes precedence over a clearly better educational judgment.

These symptoms suggest that the organization standardized \textbf{decisions requiring context} rather than standardizing the structure that helps people make those decisions well. The solution may be a stronger escalation policy, differentiated pathways, or better expert training: not abandonment of operational discipline.

\subsection*{103. How to Know When Customization Has Gone Too Far}

Watch for incompatible deadlines, unbounded office hours, a new assignment for every learner, large amounts of 1-off content, senior faculty performing routine setup, and partner-specific changes that cannot be supported later. If every sale produces a distinct product, the firm may be a bespoke service agency regardless of what its website calls it.

That can still be a sound business if priced and staffed accordingly. The danger is claiming scalable course economics while operating as an underpriced consultancy. Either productize the common job, explicitly charge for the custom portion, or accept the economics of true bespoke service.

\subsection*{104. The Stage-Gate from Repeatable Service to Institution}

Before expanding from a small founder-led service, seek evidence for 5 separate gates:

\paragraph*{1. Demand gate:} comparable buyers make real payments under honest, non-deceptive terms.

\paragraph*{2. Delivery gate:} the promised learning or operational experience can be reproduced with acceptable reliability across multiple weeks.

\paragraph*{3. Outcome gate:} meaningful customer results and independent capability gains are verified under transparent standards.

\paragraph*{4. Economics gate:} normalized contribution and cash timing can sustainably support ongoing obligations and essential reinvestment.

\paragraph*{5. Delegation gate:} another competent person or team can perform important work without the founder silently repairing every defect.

These are proposed management tests, not a universal stage model. The threshold at each gate depends on mission and risk. Research-intensive work may require long subsidized development periods; a commodity product may need much higher volume. But a service firm that crosses none of these gates should be cautious about taking on large fixed commitments.

\subsection*{105. The Institutional Production Function}

To think about the whole system without pretending to know its precise mathematical form, write:

\begin{equation}
Y = f(H, K, A, P, M; \theta),
\end{equation}
where $Y$ is the vector of meaningful customer outcomes; $H$, qualified human attention; $K$, tools and physical or digital capital; $A$, accumulated intellectual and process assets; $P$, customer prerequisites and participation; and $M$, management and coordination. $\theta$ represents environmental conditions such as market maturity, technology stability, and task complexity.

This is a conceptual production function, not an estimated econometric model. Its point is that more of 1 input can be useless when another is missing. More videos cannot compensate for no practice. More mentors cannot fully repair invalid outcomes. A new GPU cluster cannot by itself create students capable of evaluating models. Better coordination can increase useful output even when the headcount is unchanged.

The founder's task is to identify \textbf{complementarities} among inputs rather than purchasing 1 attractive asset and hoping the other pieces appear. Research on organizational complementarity warns that identifying causal synergies empirically is difficult; simultaneous changes and selection can produce misleading comparisons [E11].

\subsection*{106. Twelve Questions to Take into Any Service Business}

When inspecting an unfamiliar firm, ask 12 systematic diagnostic questions:

\begin{enumerate}
\item What is the real, unvarnished unit of value production?
\item Which activities are truly variable with customer volume?
\item What scarce resource represents the immediate binding constraint?
\item Where does work accumulate, wait in queues, and create delay?
\item Which quality failures generate the largest cascades of rework?
\item What specific tasks have become demonstrably cheaper through accumulated repetition?
\item Which intellectual, diagnostic, or operational assets are systematically reused?
\item Exactly what effort, preparation, and behavior must the customer contribute?
\item Which workflow segments can technology enhance without eroding accountability?
\item What true costs remain hidden in unpaid founder labor or deferred service liabilities?
\item How does the marginal and average cost curve shift at the next staffing threshold?
\item Is volume growth improving the customer's verified capability or merely increasing recorded units delivered?
\end{enumerate}

These questions travel well across tutoring, engineering consulting, software operations, medical practices, consumer services, laboratories, and industrial businesses. Their answers will differ; the discipline of tracing the whole system will not.

\subsection*{107. A Final Synthesis: Cheap Delivery Is Not the Purpose}

The deepest entrepreneurial mistake in service economics is to treat the human being as an expensive obstacle to scale. Sometimes human labor really is used wastefully: explaining the same basic concept repeatedly, manually formatting submissions, correcting avoidable setup failures, or transferring information between incompatible systems. These activities deserve redesign. But some human effort is the product itself: expert judgment, care in ambiguous situations, credible review, ethical accountability, and the creation of a serious community of practice.

A strong institution does not maximize humans or minimize humans. It decides \textbf{which human attention changes outcomes}, protects that attention, develops other people capable of exercising it, and uses technology and process design to remove the work that distracts from it.

A new institution may begin with the founder doing nearly everything. An enduring institution must eventually become a system in which knowledge, standards, tools, culture, and competent people reproduce the result without continuous founder heroics. The financial reward of that transition is real, but it is not the sole point. The transition is what allows the institution to serve more people without selling an increasingly hollow version of its own promise.

\begin{quote}
\itshape
The right form of educational operating leverage is the ability to make more people demonstrably capable with each unit of scarce institutional effort: while preserving the standards by which that capability is recognized.
\end{quote}
